\section{线性模型也是正交投影}

对 \(n\) 个响应值，设计矩阵的每一列描述一种允许的均值变化。最小二乘就是找列空间中离观测向量最近的点；这与前面正态样本均值方向的投影是同一种运算。

\begin{definition}[线性模型与可识别性]
设 \(X\in\mathbb R^{n\times p}\) 为已知且满列秩的设计矩阵，\(\beta\in\mathbb R^p\) 为固定未知参数，\(Y=X\beta+\varepsilon\)，其中
\(\mathbb E\varepsilon=0\)、\(\operatorname{Cov}(\varepsilon)=\sigma^2I_n\)，\(\sigma^2>0\)。
满列秩使不同 \(\beta\) 给出不同均值 \(X\beta\)，因而参数可识别。若两列相同，就只能识别它们对应系数的和，不能分别估计两系数。
\end{definition}

最小化 \(\|y-Xb\|^2\)。对 \(b\) 求导得正规方程
\(X^\top Xb=X^\top y\)。满列秩使 \(X^\top X\) 正定且可逆，因此
\[
\widehat\beta=(X^\top X)^{-1}X^\top Y,\qquad
H=X(X^\top X)^{-1}X^\top,\qquad
\widehat Y=HY,\quad e=(I-H)Y.
\]
直接计算 \(H^\top=H\)、\(H^2=H\)，故 \(H\) 是到
\(\operatorname{col}(X)\) 的正交投影，残差 \(e\) 与该列空间正交。

\begin{theorem}[Gauss--Markov，有限样本]
在上述仅有均值与协方差条件的线性模型中，任意线性无偏估计
\(\widetilde\beta=AY\)、\(AX=I_p\)，都满足
\(\operatorname{Cov}(\widetilde\beta)-\operatorname{Cov}(\widehat\beta)\) 半正定。结论不要求误差正态。
\end{theorem}
\begin{proof}
写 \(A=(X^\top X)^{-1}X^\top+D\)。由 \(AX=I_p\) 得 \(DX=0\)。
因此
\[
AA^\top=(X^\top X)^{-1}
 +(X^\top X)^{-1}X^\top D^\top
 +DX(X^\top X)^{-1}+DD^\top
 =(X^\top X)^{-1}+DD^\top .
\]
两估计协方差差为 \(\sigma^2DD^\top\)，其二次型
\(v^\top DD^\top v=\|D^\top v\|^2\ge0\)，所以半正定。
\end{proof}

若 \(X\) 的列空间包含常数向量 \(\mathbf1\)，则
\(\operatorname{span}\{\mathbf1\}\subset\operatorname{col}(X)\)，三个互相正交的部分给
\[
\|Y-\overline Y\mathbf1\|^2
=\|HY-\overline Y\mathbf1\|^2+\|(I-H)Y\|^2.
\]
这就是回归和方差分析共用的平方和分解；解释变量是连续量还是组别指示变量，只改变 \(X\) 的列，不改变投影证明。

\begin{theorem}[正态误差下的精确平方和分布]
若额外假定 \(\varepsilon\sim N_n(0,\sigma^2I_n)\)，且
\(\operatorname{rank}(X)=p<n\)，则
\(\|(I-H)Y\|^2/\sigma^2\sim\chi^2_{n-p}\)，并与 \(\widehat\beta\) 独立。若两个嵌套设计矩阵 \(X_0,X_1\) 的列空间满足
\(\operatorname{col}(X_0)\subset\operatorname{col}(X_1)\)，秩分别为
\(p_0<p_1\)，且零假设中的均值属于 \(\operatorname{col}(X_0)\)，则
\[
F=\frac{\{\mathrm{RSS}_0-\mathrm{RSS}_1\}/(p_1-p_0)}
{\mathrm{RSS}_1/(n-p_1)}
\sim F_{p_1-p_0,n-p_1}.
\]
\end{theorem}
\begin{proof}
因 \((I-H)X\beta=0\)，残差为 \((I-H)\varepsilon\)；
\(I-H\) 是秩 \(n-p\) 的正交投影。上一章的高斯投影定理给第一结论，也给它与 \(H\varepsilon\) 独立。对嵌套模型，\(H_1-H_0\) 和 \(I-H_1\) 是像互相正交、秩分别为 \(p_1-p_0\)、\(n-p_1\) 的投影。在零假设下它们也都消去均值。
\(\mathrm{RSS}_0-\mathrm{RSS}_1=Y^\top(H_1-H_0)Y\)，
\(\mathrm{RSS}_1=Y^\top(I-H_1)Y\)；
两者除以 \(\sigma^2\) 是独立卡方，按 \(F\) 分布的定义得式。
\end{proof}

同一个投影证明也给出回归系数的精确区间。取一个事先指定的非零向量 \(c\in\mathbb R^p\)，要估计线性目标 \(c^\top\beta\)。由正规方程可直接写
\(\widehat\beta-\beta=(X^\top X)^{-1}X^\top\varepsilon\)，所以
\[
\mathbb E(c^\top\widehat\beta)=c^\top\beta,
\qquad
\operatorname{Var}(c^\top\widehat\beta)
=\sigma^2c^\top(X^\top X)^{-1}c.
\]
残差平方和的期望为 \((n-p)\sigma^2\)，故
\(s^2=\mathrm{RSS}/(n-p)\) 无偏。在正态误差条件下，标准化的
\(c^\top\widehat\beta-c^\top\beta\) 是标准正态，且因为它只依赖
\(H\varepsilon\)，与只依赖 \((I-H)\varepsilon\) 的 \(s^2\) 独立。于是
\[
\frac{c^\top\widehat\beta-c^\top\beta}
{s\sqrt{c^\top(X^\top X)^{-1}c}}
\sim t_{n-p}.
\]
从而把绝对值限制在 \(t_{n-p,1-\alpha/2}\) 内，便得到 \(c^\top\beta\) 的精确 \(1-\alpha\) 置信区间。若误差只有相同方差但不正态，Gauss--Markov 的最优线性无偏结论仍成立，这个有限样本 \(t\) 分布却不再有保证。

投影矩阵的对角元 \(h_{ii}=e_i^\top He_i\) 衡量第 \(i\) 个观测位置对自身拟合值的影响，称为杠杆值。因 \(H\) 是正交投影，
\(0\le h_{ii}=\|He_i\|^2\le1\)，且
\(\sum_i h_{ii}=\operatorname{tr}H=p\)。残差协方差为
\(\sigma^2(I-H)\)，故
\(\operatorname{Var}(e_i)=\sigma^2(1-h_{ii})\)。高杠杆位置的残差可能偏小，不等于该观测可靠；它也可能把拟合直线拉向自己。这个结论直接来自投影，不能仅凭残差绝对值判断异常观测。

最后，若已知 \(\operatorname{Cov}(\varepsilon)=\Sigma\) 为正定而非
\(\sigma^2I\)，以 \(\Sigma^{-1/2}\) 同时变换 \(Y\) 和 \(X\)，便把误差白化。对变换后的问题应用同一投影证明，得到广义最小二乘
\[
\widehat\beta_{\mathrm{GLS}}
=(X^\top\Sigma^{-1}X)^{-1}X^\top\Sigma^{-1}Y.
\]
\(\Sigma\) 若由数据估计，上面基于已知协方差和正态投影的有限样本分布一般不再精确。

满列秩的假设值得单独检查一次，因为它同时承担了可识别性与稳定性两件事。若两列几乎相同，\(X^\top X\) 的最小本征值 \(s_{\min}^2\) 很小，而 \(\widehat\beta\) 的协方差含 \((X^\top X)^{-1}\)，噪声沿该方向被 \(1/s_{\min}^2\) 放大。这与\cref{sec:math-ill-posed-inverse-problem}中小奇异值放大数据误差是同一个现象，只是这里的方向来自设计矩阵而不是积分核。一个直接的补救是把正规方程改成
\[
\widehat\beta_\alpha=(X^\top X+\alpha I)^{-1}X^\top Y .
\]
沿 \(X\) 的右奇异向量 \(v_i\) 分解，它把每个系数按
\(\lambda_i/(\lambda_i+\alpha)\) 收缩（\(\lambda_i\) 是 \(X^\top X\) 的本征值），\cref{fig:math-ridge-spectrum}画出这些因子。\(\lambda_i\gg\alpha\) 的方向几乎不动，\(\lambda_i\ll\alpha\) 的方向被压向零。代价是 \(\widehat\beta_\alpha\) 有偏：它不再是无偏估计，Gauss--Markov 的最优性也就不再适用。用偏差换方差是否划算，取决于 \(\sigma^2\) 与 \(\alpha\) 的相对大小，没有普遍答案；这里能确定的是，滤波因子 \(\lambda/(\lambda+\alpha)\) 与前一章 Tikhonov 的 \(s/(s^2+\alpha)\) 形状完全相同。

\begin{figure}[htbp]
\centering
\includegraphics[width=.76\linewidth]{ridge_spectrum.pdf}
\caption{岭回归谱收缩 $\lambda_i/(\lambda_i+\alpha)$。不同 $\alpha$ 下，小特征值被压向零、大特征值保持不变，虚线为无收缩情形。}
\label{fig:math-ridge-spectrum}
\end{figure}

\begin{exercise}
取两组各两次观测，写出包含截距与组别指示变量的 \(X\)，求其投影矩阵。再将两组均值相等写为嵌套模型检验，并验证残差自由度为 \(2\)。
\end{exercise}
